\clearpage
\FrontTitle{ABSTRAK}
\addcontentsline{toc}{chapter}{ABSTRAK}
\begin{center}
{\bfseries \JudulLaporan\par}
\vspace{0.3cm}
Oleh\par
\NamaRaif{} (\NIMRaif)\\
\NamaFrederiko{} (\NIMFrederiko)
\end{center}
\vspace{0.3cm}
Sistem pengelolaan alur kerja memerlukan kesesuaian antara penyajian informasi, pengaturan hak akses, dan pencatatan perkembangan pekerjaan. Kerja praktik di PT Akhdani Reka Solusi pada \PeriodeKP{} dilaksanakan untuk berkontribusi pada pengembangan dan penjaminan kualitas bagian sistem tersebut. Tahapan pekerjaan meliputi pemahaman kebutuhan dan kode yang telah tersedia, perancangan perubahan, implementasi antarmuka dan layanan, integrasi, serta penyusunan pengujian. Pekerjaan pengembangan menghasilkan halaman dasbor dengan layanan data simulasi, fitur pengelolaan peran beserta hak akses, dan riwayat alur kerja yang menampilkan aktivitas serta transisi. Pengujian komponen disusun untuk memeriksa pemuatan data, interaksi, validasi, dan penanganan kegagalan layanan. Perubahan kebutuhan pada pertengahan pelaksanaan diikuti pembelajaran Go dan adaptasi fitur pengelolaan peran ke layanan baru. Pada akhir periode juga dihasilkan rancangan layar pengelolaan dasbor menggunakan Visily. \IsianFrederiko{ringkasan metode dan hasil penjaminan kualitas}. Pelaksanaan kerja praktik menunjukkan pentingnya keselarasan data antarlapisan aplikasi, pemeriksaan perilaku pengguna, dan kemampuan menyesuaikan implementasi ketika teknologi berubah. Kegiatan ini turut memberikan pengalaman koordinasi dan pemeliharaan kode dalam tim pengembangan perangkat lunak.

\vspace{0.3cm}
\noindent\textbf{Kata kunci:} alur kerja, pengembangan perangkat lunak, hak akses, pengujian perangkat lunak, kerja praktik.
